% This note was written entirely by Claude (Fable 5, via Claude Code), from
% the derivations to the numerical verification against the cached Vern9
% references; see the commit message for provenance.
\documentclass[11pt]{article}

\usepackage[margin=1.1in]{geometry}
\usepackage{amsmath, amssymb, mathtools}
\usepackage{bm}
\usepackage{booktabs}
\usepackage[colorlinks=true, linkcolor=blue, citecolor=blue]{hyperref}

\newcommand{\ii}{\mathrm{i}}
\renewcommand{\Re}{\operatorname{Re}}
\renewcommand{\Im}{\operatorname{Im}}
\newcommand{\code}[1]{\texttt{#1}}
\newcommand{\tr}{\operatorname{tr}}
\newcommand{\norm}[1]{\lVert #1 \rVert}
\DeclarePairedDelimiter{\abs}{\lvert}{\rvert}

\title{Why a $6\times10^{-3}$ RWA error moves the CNOT gate fidelity\\
by only $2\times10^{-6}$}
\author{}
\date{\today}

\begin{document}
\maketitle

\section{The apparent paradox}

The CNOT3 convergence experiment solves the same physical gate problem in the
rotating frame with the rotating wave approximation (\code{rwa}) and without it
(\code{norwa}, exactly equivalent to the laboratory frame).  The convergence
figures mark the ``RWA modeling error,'' the 2-norm difference between the two
converged final states,
\begin{equation}
    \norm{\delta} \coloneqq \norm{U_{\mathrm{norwa}} - U_{\mathrm{rwa}}}_F
    = 6.008\times10^{-3},
\end{equation}
as the accuracy floor below which further numerical refinement of the RWA model
is pointless.  Yet the CNOT gate fidelities of the two solutions, measured
against the same target, are nearly identical:
\begin{equation}
    1 - F_{\mathrm{rwa}} = 4.0269\times10^{-5},
    \qquad
    1 - F_{\mathrm{norwa}} = 4.2667\times10^{-5},
\end{equation}
a difference of only $2.398\times10^{-6}$, some $2500$ times smaller than
$\norm{\delta}$.  Both numbers are correct; they measure different things, and
the gap between them is the product of two geometric effects.  This note works
them out and verifies each one quantitatively against the computed solutions.

\paragraph{Setup.}  The system has $N = 160$ levels (a 10-level storage cavity
and two 4-level transmons) and $m = 4$ essential computational states.  A
solution is the matrix $U \in \mathbb{C}^{N\times m}$ whose columns are the
essential basis states propagated to the final time $T = 550$~ns; unitarity
gives $\norm{U}_F = \sqrt{m}$.  The target is the isometry
$V \in \mathbb{C}^{N\times m}$ (satisfying $V^\dagger V = I_m$) whose columns
are the CNOT images of the basis states, expressed in the appropriate frame.
The gate fidelity is the standard trace overlap
\begin{equation} \label{eq:fidelity}
    F(U) = \frac{\abs{z}^2}{m^2},
    \qquad
    z \coloneqq \tr\bigl(V^\dagger U\bigr) = \langle V, U\rangle_F .
\end{equation}
Because the columns of $V$ are supported only on the essential levels, the
overlap automatically projects out guard-level leakage; leakage lowers
$\abs{z}$ and hence $F$.

\section{Effect 1: infidelity is quadratic in the miss distance}

Fidelity is blind to a global phase, so the natural distance from a solution to
the target is the distance to the whole phase orbit
$\{e^{\ii\varphi} V : \varphi \in [0, 2\pi)\}$:
\begin{equation}
    d(U) \coloneqq \min_{\varphi} \,\norm{U - e^{\ii\varphi} V}_F .
\end{equation}
Expanding the square and minimizing (the optimum is
$e^{\ii\varphi^*} = z/\abs{z}$) gives
$d^2 = \norm{U}_F^2 + \norm{V}_F^2 - 2\abs{z} = 2m - 2\abs{z}$, i.e.\
$\abs{z} = m - d^2/2$.  Substituting into \eqref{eq:fidelity},
\begin{equation} \label{eq:quadratic}
    \boxed{\;
    1 - F = \frac{d^2}{m}\Bigl(1 - \frac{d^2}{4m}\Bigr)
          \approx \frac{d^2}{m}
    \quad (d \ll 1). \;}
\end{equation}
Infidelity is \emph{quadratic} in how far the solution sits from the target.
A miss of size $10^{-2}$ therefore lives on the $10^{-5}$ infidelity scale,
with no error mechanism required beyond the change of units.  Numerically,
\begin{equation}
    d_{\mathrm{rwa}} = 1.2692\times10^{-2}
    \;\Rightarrow\;
    \frac{d_{\mathrm{rwa}}^2}{4} = 4.0270\times10^{-5},
\end{equation}
which reproduces the measured $1 - F_{\mathrm{rwa}} = 4.0269\times10^{-5}$ to
four digits (the small deficit is the $d^4$ term).  Likewise
$d_{\mathrm{norwa}} = 1.3064\times10^{-2}$ gives $4.267\times10^{-5}$.

Effect 1 alone does not resolve the paradox: if the RWA perturbation $\delta$
pushed the solution straight away from the target, the infidelity would change
by up to $(2 d\,\norm{\delta} + \norm{\delta}^2)/m = 4.7\times10^{-5}$, i.e.\ it
would roughly double.  The observed change is twenty times smaller than that
bound, which is the second effect.

\section{Effect 2: fidelity sees only the radial component}

The infidelity \eqref{eq:quadratic} depends on the state only through the
scalar $d$, the distance to the target orbit.  Under a perturbation
$U \mapsto U + \delta$, the overlap shifts by
$\langle V, \delta\rangle$ and, to first order,
\begin{equation}
    \Delta\abs{z}
    = \Re\Bigl[\tfrac{\bar z}{\abs{z}}\,\langle V, \delta\rangle\Bigr]
    + O\bigl(\norm{\delta}^2\bigr),
\end{equation}
the projection of $\delta$ onto a \emph{single real direction} out of the
$2Nm = 1280$ real dimensions of the state space: the ``radial'' direction
$\hat r = (U - e^{\ii\varphi^*}V)/d$ pointing from the phase-aligned target to
the solution.  Every other direction is tangential and invisible at first
order.  (Since $d$ is a distance function, it is also globally 1-Lipschitz,
$\abs{\Delta d} \le \norm{\delta}$, so radial alignment is the worst case, not
just the first-order one.)

Decomposing the measured RWA perturbation $\delta$ at $U_{\mathrm{rwa}}$ into
mutually orthogonal pieces gives Table~\ref{tab:decomposition}.

\begin{table}[ht]
\centering
\begin{tabular}{llrr}
\toprule
component & direction & size & share of $\norm{\delta}^2$ \\
\midrule
global phase & $\ii U/\norm{U}_F$ & $5.211\times10^{-3}$ & $75.2\%$ \\
other tangential & (orthogonal remainder) & $2.990\times10^{-3}$ & $24.8\%$ \\
radial & $\hat r$ & $2.58\times10^{-5}$ & $0.002\%$ \\
\midrule
total & & $6.008\times10^{-3}$ & $100\%$ \\
\bottomrule
\end{tabular}
\caption{Orthogonal decomposition of the RWA perturbation
$\delta = U_{\mathrm{norwa}} - U_{\mathrm{rwa}}$.}
\label{tab:decomposition}
\end{table}

Three-quarters of $\norm{\delta}^2$ is a \emph{relative global phase}: the two
solutions differ by $e^{\ii\Delta\varphi}$ with $\Delta\varphi = 2.6$~mrad,
equivalent to a common level shift of under a kilohertz integrated over the
550~ns gate.  Fidelity is exactly invariant under this component, at all
orders.  Almost all of the rest is tangential: it changes the \emph{direction}
in which the solution misses the target, not the distance.  A tangential step
of length $\tau$ does grow the distance at second order, through the curvature
of the sphere of radius $d$ around the target orbit,
$\Delta d \approx \tau^2/(2d)$, and this closes the books exactly:
\begin{equation}
    \underbrace{\frac{(2.990\times10^{-3})^2}{2\,(1.2692\times10^{-2})}}_{\text{tangential curvature}}
    + \underbrace{2.58\times10^{-5}}_{\text{radial}}
    = 3.78\times10^{-4}
    \;\approx\;
    \Delta d_{\mathrm{measured}} = 3.72\times10^{-4},
\end{equation}
and then, by \eqref{eq:quadratic},
\begin{equation}
    \Delta(1-F)
    = \frac{d_{\mathrm{norwa}}^2 - d_{\mathrm{rwa}}^2}{m}
    = 2.398\times10^{-6},
\end{equation}
which is the measured infidelity difference.

The tiny radial component is not a coincidence of tuning.  A perturbation
pointing in a \emph{random} direction of the 1280-dimensional real state space
would have an expected radial fraction of order $1/\sqrt{1280} \approx 2.8\%$;
the measured first-order radial fraction is $0.43\%$.  The counter-rotating
terms have no knowledge of the particular direction in which the optimized
pulse happens to miss the target, so their effect is generically tangential.
Physically this is the expected structure of the RWA error: the leading effect
of the counter-rotating terms is a small coherent energy shift (of
Bloch--Siegert type~\cite{BlochSiegert1940, Zeuch2020}), whose dominant
consequence over a long gate is accumulated phase.  The common part of that
phase is pure gauge (the $75\%$ row of Table~\ref{tab:decomposition}), and the
differential part redistributes the miss rather than enlarging it.

\section{Reconciling the two plots}

Both numbers are the right answer to different questions.  The 2-norm line
$\norm{\delta} = 6\times10^{-3}$ answers: \emph{below what numerical error does
refining the RWA model stop improving agreement with the true (lab-frame)
state?}  For full-state observables, phases included, the RWA really is a
$6\times10^{-3}$-level modeling error.  The fidelity comparison answers:
\emph{does the RWA change how good the gate is?}  There the two effects stack:
the quadratic map \eqref{eq:quadratic} converts $10^{-2}$-scale distances to
$10^{-5}$-scale infidelities, and the radial projection suppresses the
sensitivity to this particular perturbation by another factor of twenty,
leaving $2.4\times10^{-6}$.  A useful mnemonic for the overall bookkeeping:
\begin{equation}
    \underbrace{6.0\times10^{-3}}_{\norm{\delta}}
    \;\xrightarrow{\text{drop global phase}}\;
    3.0\times10^{-3}
    \;\xrightarrow{\text{curvature } \tau^2/2d}\;
    \Delta d = 3.7\times10^{-4}
    \;\xrightarrow{\;2d\,\Delta d/m\;}\;
    2.4\times10^{-6}.
\end{equation}

\section{Relation to the literature}

The two ingredients are individually well documented, though we are not aware
of a reference that presents this specific side-by-side observation (state-norm
RWA error versus gate-fidelity RWA effect) for an optimized multilevel gate.

\emph{Quadratic gap for coherent errors.}  Effect 1 is the same mathematics
behind a well-known result in the gate-benchmarking literature: for coherent
(unitary) errors the worst-case error rate (diamond norm) scales as the
\emph{square root} of the average gate infidelity, so norm-type error measures
and fidelity-type measures of the same coherent error can differ by orders of
magnitude~\cite{Sanders2016, Kueng2016, Wallman2016}.  That literature runs the
argument in the direction ``reported fidelities flatter coherent errors'';
here the same fact appears in reverse, as a plotted norm error that looks
alarming on a fidelity scale.

\emph{Stationarity of fidelity.}  Effect 2 is a consequence of $F$ being
stationary in every direction but one; this is the elementary geometry
underlying the quantum control landscape literature, where the fidelity's
critical structure as a function of the evolution is analyzed in
detail~\cite{Rabitz2004}.

\emph{Character of the RWA error.}  That the leading counter-rotating effect
is a coherent, largely calibratable phase shift goes back to Bloch and
Siegert~\cite{BlochSiegert1940}; modern treatments include the exact rotating
frame construction of Zeuch et al.~\cite{Zeuch2020} and recent gate-focused
analyses that compute counter-rotating contributions to infidelity
perturbatively and cancel them with pulse corrections or phase
calibration~\cite{Rower2024, Fluxonium2025}.  Those works are consistent with
Table~\ref{tab:decomposition}: the RWA error is dominated by phase
accumulation, which gate fidelity largely forgives.

\begin{thebibliography}{9}

\bibitem{Sanders2016}
Y.~R. Sanders, J.~J. Wallman, and B.~C. Sanders,
\emph{Bounding quantum gate error rate based on reported average fidelity},
New J. Phys. \textbf{18}, 012002 (2016), arXiv:1501.04932.

\bibitem{Kueng2016}
R.~Kueng, D.~M. Long, A.~C. Doherty, and S.~T. Flammia,
\emph{Comparing experiments to the fault-tolerance threshold},
Phys. Rev. Lett. \textbf{117}, 170502 (2016), arXiv:1510.05653.

\bibitem{Wallman2016}
J.~J. Wallman,
\emph{Bounding experimental quantum error rates relative to fault-tolerant
thresholds}, arXiv:1511.00727 (2016).

\bibitem{Rabitz2004}
H.~Rabitz, M.~Hsieh, and C.~Rosenthal,
\emph{Quantum optimally controlled transition landscapes},
Science \textbf{303}, 1998 (2004).

\bibitem{BlochSiegert1940}
F.~Bloch and A.~Siegert,
\emph{Magnetic resonance for nonrotating fields},
Phys. Rev. \textbf{57}, 522 (1940).

\bibitem{Zeuch2020}
D.~Zeuch, F.~Hassler, J.~J. Slim, and D.~P. DiVincenzo,
\emph{Exact rotating wave approximation},
Ann. Phys. \textbf{423}, 168327 (2020), arXiv:1807.02858.

\bibitem{Rower2024}
D.~A. Rower et al.,
\emph{Suppressing counter-rotating errors for fast single-qubit gates with
fluxonium}, PRX Quantum \textbf{5}, 040342 (2024).

\bibitem{Fluxonium2025}
\emph{Single-qubit gates beyond the rotating-wave approximation for strongly
anharmonic low-frequency qubits}, arXiv:2503.08238 (2025).

\end{thebibliography}

\appendix

\section{Reproducing the numbers}

All quantities in this note come from the cached Vern9 references
($\code{abstol} = \code{reltol} = 10^{-15}$, basis initial condition,
$T = 550$, $N_{\mathrm{osc}} = 10$, two guard levels per transmon) via
\code{vern9\_reference} and \code{cnot3\_gate\_fidelity} in
\code{src/cnot3\_reference.jl}.  The decomposition of
Table~\ref{tab:decomposition} projects
$\delta = U_{\mathrm{norwa}} - U_{\mathrm{rwa}}$ onto the gauge direction
$\ii U_{\mathrm{rwa}}/\norm{U_{\mathrm{rwa}}}_F$ and the (gauge-orthogonalized)
radial direction $(U_{\mathrm{rwa}} - e^{\ii\varphi^*}V)/d$, using the real
inner product $\Re\langle\cdot,\cdot\rangle_F$; the tangential remainder is
obtained by Pythagoras.

\end{document}
